\PassOptionsToPackage{usenames,dvipsnames}{xcolor}
\documentclass[a4paper,12pt]{letter}
\usepackage[utf8]{inputenc}
\usepackage[T1]{fontenc}
\usepackage[english,slovene]{babel}
\usepackage[tikz]{bclogo}
\usepackage{qrcode,marvosym}
\usepackage{pgf,tikz,pgfplots}
\pgfplotsset{compat=1.14}
\usepackage{mathrsfs}
\usetikzlibrary{arrows}
\usepackage{graphicx}
\newcommand{\ds}{\displaystyle}
\usepackage{fancybox}
\usepackage{amsfonts}
\usepackage{amsmath}
\usepackage{wallpaper}
\usepackage{hyperref}
\usepackage{array}
\newcolumntype{M}[1]{>{\centering\arraybackslash}m{#1}}
\newcolumntype{N}{@{}m{0pt}@{}}
\usepackage{fancyhdr,enumerate}
\newcommand{\ora}{\overrightarrow}
\renewcommand{\headrulewidth}{3pt}
\renewcommand{\headrule}{\hbox to\headwidth{%
  \color{gray}\leaders\hrule height \headrulewidth\hfill}}
\renewcommand{\footrulewidth}{2pt}
\renewcommand{\footrule}{\hbox to\headwidth{%
  \color{brown}\leaders\hrule height \footrulewidth\hfill}}
\usepackage{gensymb}
\usepackage{amssymb}
  \textwidth 20 true cm
  \oddsidemargin -2 true cm
  \evensidemargin -2 true cm
  \topmargin -3.5 true cm
  \textheight 27.5 true cm
  \linespread{1.2}
\renewcommand{\headrulewidth}{0.3pt}
\renewcommand{\footrulewidth}{1.9pt}
\definecolor{qqwuqq}{rgb}{1,0,0}
\definecolor{qqqqff}{rgb}{0,0,1}
\definecolor{qqffqq}{rgb}{0,1,0}
\definecolor{ffqqqq}{rgb}{1,0,0}
\definecolor{zzuxux}{rgb}{0,0,0}
\usepackage{mdframed}
\usepackage{multicol}
\definecolor{bgblue}{RGB}{0,0,253}
\usepackage{eurosym}
%%%%%%%%%%%%%%%%%%%%%%%%%%%%%%%%%%%%%%%%%%%%%%%%%%%%%%%%
\newcounter{tocke}
\setcounter{tocke}{0}
\usepackage{tikz-3dplot}
%%%%%%%%%%%%%%%%%%%%%%%%%%%%%%%%%%%%%%%%%%%%%%%%%%%%%%%%%%
\mdfdefinestyle{theoremstyle}{%
linecolor=brown!40,linewidth=1pt,%
frametitlerule=true,%
frametitlebackgroundcolor=brown!50,
innertopmargin=\topskip,
}
\mdfdefinestyle{theoremstyle1}{%
linecolor=brown,middlelinewidth=1pt,%
frametitlerule=true,%
apptotikzsetting={\tikzset{mdfframetitlebackground/.append style={%
shade,left color=black!40, right color=gray!10},
mdfbackground/.append style={
top color=white!100!white,
bottom color=black!5
},
}},
frametitlerulecolor=gray,
frametitlerulewidth=2pt,
innertopmargin=\topskip,
innerbottommargin=\topskip,
}
\mdtheorem[style=theoremstyle1]{naloga}{Naloga}
\mdtheorem[style=theoremstyle]{resitev}{Analiza Naloge}
\mdtheorem[style=theoremstyle]{kriterij}{\v{S}tevilo dose\v{z}enih to\v{c}k na testu}
%%%%%%%%%%%%%%%%%%%%%%%%%%%%%%%%%%%%%%%%%%%%%%%%%%%%%%%%%%
\usepackage[table]{xcolor}
\usepackage[printwatermark]{xwatermark}
\newwatermark[allpages,color=brown!2,angle=45,scale=5,
xpos=-5,ypos=0]{matej.info}
%%%%%%%%%%%%%%%%%%%%%%%%%%%%%%%%%%%%%%%%%%%%%%%%%%%%%%%%%%

\begin{document}
\pagestyle{fancy}
\renewcommand\headrulewidth{0pt}
\lhead{}\chead{}\rhead{}
\rfoot{\vspace*{-2\baselineskip}}
\renewcommand{\vec}{\overrightarrow}
%%%%%%%%%%%%%%%%%%%% podatki testa -- tu spreminjaj
\def\programletnik{Gim - 4}
\def\naslov{Test 1.2 - }
\def\sklop{Geometrijska telesa}
%%%%%%%%%%%%%%%%%%%% tocke po nalogah
\def\tockeA{4+6}
\def\tockeB{6+3}
\def\tockeC{5+4+3}
\def\tockeD{5+4}
%%%%%%%%%%%%%%%%%%%%
\begin{minipage}{20cm}
\begin{bclogo}[couleur=brown!30,
  arrondi=0,ombre=true, couleurOmbre=gray!50,
logo =\bcplume,
  noborder=true]
{\textrm{{\Large{\naslov \sklop} \Huge{}
\large{} \hfill $\programletnik$}}}
\end{bclogo}
\end{minipage}
\begin{minipage}{20cm}
\begin{bclogo}[couleur=brown!30,
  arrondi=0,ombre=true, couleurOmbre=gray!50,
logo =\bcplume,
  noborder=true]
{\textrm{{\raisebox{-0.3cm}{\large{Ime in Priimek: \rule{10cm}{0.1mm}} \Huge{}
\large{} \hfill $ $}}}}
\end{bclogo}
\end{minipage}
%%%%%%%%%%%%%%%%%%%%%%%%%%%%%%%%%%%%%%%%%%%%%%%%%

%%%%%%%%%%%%%%%%%%%%%%%%%%%%%%%%%%%%%%%%%%%%%%%%%
\begin{naloga}[\hfill $\tockeA$] \addtocounter{tocke}{\numexpr\tockeA\relax}
Polkroglo s prostornino $V=144\pi\ \text{cm}^3$ pretalimo v enakostrani\v{c}ni valj (osni presek valja je kvadrat).
\begin{enumerate}[a)]
\item Izra\v{c}unaj povr\v{s}ino polkrogle (skupaj z osnovno ploskvijo).
\item Izra\v{c}unaj povr\v{s}ino pla\v{s}\v{c}a valja.
\end{enumerate}
\end{naloga}\newpage

%%%%%%%%%%%%%%%%%%%%%%%%%%%%%%%%%%%%%%%%%%%%%%%%%
\begin{naloga}[\hfill $\tockeB$] \addtocounter{tocke}{\numexpr\tockeB\relax}
Pravilni tristrani prizmi meri povr\v{s}ina pla\v{s}\v{c}a $S_{pl}=90\ \text{cm}^2$, plo\v{s}\v{c}ina osnovne ploskve pa je $S_o=9\sqrt3\ \text{cm}^2$.
\begin{enumerate}[a)]
\item Izra\v{c}unaj prostornino prizme.
\item Izra\v{c}unaj povr\v{s}ino prizme.
\end{enumerate}
\end{naloga}\newpage

%%%%%%%%%%%%%%%%%%%%%%%%%%%%%%%%%%%%%%%%%%%%%%%%%
\begin{naloga}[\hfill $\tockeC$] \addtocounter{tocke}{\numexpr\tockeC\relax}
Pravilni \v{s}tiristrani piramidi meri povr\v{s}ina pla\v{s}\v{c}a $S_{pl}=48\ \text{cm}^2$, vi\v{s}ina stranske ploskve (stranska vi\v{s}ina) pa $m=4$ cm.
\begin{enumerate}[a)]
\item Izra\v{c}unaj prostornino piramide.
\item Izra\v{c}unaj vsoto dol\v{z}in vseh robov piramide.
\item Izra\v{c}unaj kot med stransko in osnovno ploskvijo (do kotne minute natan\v{c}no).
\end{enumerate}
\end{naloga}\newpage

%%%%%%%%%%%%%%%%%%%%%%%%%%%%%%%%%%%%%%%%%%%%%%%%%
\begin{naloga}[\hfill $\tockeD$] \addtocounter{tocke}{\numexpr\tockeD\relax}
Plo\v{s}\v{c}ina osnega preseka pokončnega sto\v{z}ca je $60\ \text{cm}^2$, vi\v{s}ina sto\v{z}ca pa je $v=12$ cm.
\begin{enumerate}[a)]
\item Izra\v{c}unaj povr\v{s}ino pla\v{s}\v{c}a sto\v{z}ca.
\item Izra\v{c}unaj polmer sto\v{z}cu včrtane krogle.
\end{enumerate}
\end{naloga}

\vfill
\begin{kriterij*} \begin{center} \renewcommand{\arraystretch}{2.5}
\begin{tabular}{|>{\columncolor{gray!20}}c||c|c|c|c|c||c|c|c|} \hline Odstotki &$[0,45)$ & $[45,60)$ & $[60,75)$ & $[75,90)$ & $[90,100]$& \v{S}t. mo\v{z}nih to\v{c}k& Dose\v{z}ene to\v{c}ke & OCENA \\ \hline Ocena & 1 & 2 & 3 & 4 & 5 &\cellcolor{gray!20} $\thetocke$&\cellcolor{gray!20} \phantom{ocena} & \cellcolor{gray!20} \phantom{a} \\ \hline \end{tabular} \end{center} \vspace{0.3cm}  \end{kriterij*}

\newpage
\textsc{\Large Re\v{s}itve s postopki}
\vspace{0.4cm}

%%%%%%%%%%%%%%%%%%%% pregled rezultatov
\begin{center}
\renewcommand{\arraystretch}{2.1}
\begin{tabular}{|>{\columncolor{gray!20}}c|c|c|c|}
\hline
\textbf{Naloga} & \textbf{a)} & \textbf{b)} & \textbf{c)} \\ \hline
\textbf{1} & $S=108\pi\ \text{cm}^2\approx339{,}29\ \text{cm}^2$ & $S_{pl}=48\pi\sqrt[3]{3}\ \text{cm}^2\approx217{,}49\ \text{cm}^2$ & -- \\ \hline
\textbf{2} & $V=45\sqrt3\ \text{cm}^3\approx77{,}94\ \text{cm}^3$ & $S=90+18\sqrt3\ \text{cm}^2\approx121{,}18\ \text{cm}^2$ & -- \\ \hline
\textbf{3} & $V=12\sqrt7\ \text{cm}^3\approx31{,}75\ \text{cm}^3$ & $o=44\ \text{cm}$ & $\varphi\approx41\degree25'$ \\ \hline
\textbf{4} & $S_{pl}=65\pi\ \text{cm}^2\approx204{,}20\ \text{cm}^2$ & $\rho=\dfrac{10}{3}\ \text{cm}\approx3{,}33\ \text{cm}$ & -- \\ \hline
\end{tabular}
\end{center}
\vspace{0.2cm}

%%%%%%%%%%%%%%%%%%%% naloga 1
\begin{resitev}[-- polkrogla in enakostrani\v{c}ni valj]
\textbf{a)} Prostornina polkrogle s polmerom $R$ je $V=\dfrac23\pi R^3$. Iz $V=144\pi$ dobimo
\[
\frac23\pi R^3=144\pi\ \Rightarrow\ R^3=216\ \Rightarrow\ R=6\ \text{cm}.
\]
Povr\v{s}ina polkrogle je vsota ukrivljene ploskve (polovica sfere) in osnovne ploskve (kroga):
\[
S=2\pi R^2+\pi R^2=3\pi R^2=3\pi\cdot36=108\pi\ \text{cm}^2.
\]
\[
\boxed{S=108\pi\ \text{cm}^2\approx339{,}29\ \text{cm}^2}
\]

\textbf{b)} Pri pretalitvi se prostornina ohrani: $V_{valj}=144\pi$. Enakostrani\v{c}ni valj ima vi\v{s}ino enako premeru osnovne ploskve, torej $h=2r$:
\[
V_{valj}=\pi r^2h=\pi r^2\cdot2r=2\pi r^3=144\pi\ \Rightarrow\ r^3=72\ \Rightarrow\ r=\sqrt[3]{72}=2\sqrt[3]{9}\ \text{cm}\approx4{,}16\ \text{cm}.
\]
Povr\v{s}ina pla\v{s}\v{c}a valja (z $h=2r$):
\[
S_{pl}=2\pi rh=4\pi r^2=4\pi\left(2\sqrt[3]{9}\right)^2=16\pi\sqrt[3]{81}=16\pi\cdot3\sqrt[3]{3}=48\pi\sqrt[3]{3}\ \text{cm}^2.
\]
\[
\boxed{S_{pl}=48\pi\sqrt[3]{3}\ \text{cm}^2\approx217{,}49\ \text{cm}^2}
\]
\end{resitev}

\newpage
%%%%%%%%%%%%%%%%%%%% naloga 2
\begin{resitev}[-- tristrana prizma]
Osnovna ploskev je enakostrani\v{c}en trikotnik z robom $a$, vi\v{s}ina prizme je $v$.

\textbf{a)} Rob osnovne ploskve dobimo iz njene plo\v{s}\v{c}ine:
\[
S_o=\frac{a^2\sqrt3}{4}=9\sqrt3\ \Rightarrow\ a^2=36\ \Rightarrow\ a=6\ \text{cm}.
\]
Pla\v{s}\v{c} sestavljajo trije skladni pravokotniki $a\times v$:
\[
S_{pl}=3av=90\ \Rightarrow\ 3\cdot6\cdot v=90\ \Rightarrow\ v=5\ \text{cm}.
\]
Prostornina prizme:
\[
V=S_o\cdot v=9\sqrt3\cdot5=45\sqrt3\ \text{cm}^3.
\]
\[
\boxed{V=45\sqrt3\ \text{cm}^3\approx77{,}94\ \text{cm}^3}
\]

\textbf{b)} Povr\v{s}ina prizme je vsota dveh osnovnih ploskev in pla\v{s}\v{c}a:
\[
S=2S_o+S_{pl}=2\cdot9\sqrt3+90=90+18\sqrt3\ \text{cm}^2.
\]
\[
\boxed{S=90+18\sqrt3\ \text{cm}^2\approx121{,}18\ \text{cm}^2}
\]
\end{resitev}

\newpage
%%%%%%%%%%%%%%%%%%%% naloga 3
\begin{resitev}[-- \v{s}tiristrana piramida]
Osnova je kvadrat z robom $a$ (diagonala $d=a\sqrt2$), vi\v{s}ina piramide je $v$, vi\v{s}ina stranske ploskve (apotema) je $m=4$ cm, stranski rob pa $b$.

\textbf{a)} Pla\v{s}\v{c} sestavljajo \v{s}tirje skladni enakokraki trikotniki z osnovnico $a$ in vi\v{s}ino $m$:
\[
S_{pl}=4\cdot\frac{am}{2}=2am=48\ \Rightarrow\ 2\cdot a\cdot4=48\ \Rightarrow\ a=6\ \text{cm}.
\]
Vi\v{s}ino piramide dobimo iz pravokotnega trikotnika s katetama $v$ in $\frac a2=3$ cm ter hipotenuzo $m=4$ cm:
\[
v=\sqrt{m^2-\left(\frac a2\right)^2}=\sqrt{16-9}=\sqrt7\ \text{cm}.
\]
Prostornina piramide:
\[
V=\frac13a^2v=\frac13\cdot36\cdot\sqrt7=12\sqrt7\ \text{cm}^3.
\]
\[
\boxed{V=12\sqrt7\ \text{cm}^3\approx31{,}75\ \text{cm}^3}
\]

\textbf{b)} Piramida ima $4$ robove osnovne ploskve (vsak dolg $a$) in $4$ stranske robove (vsak dolg $b$). Stranski rob dobimo iz pravokotnega trikotnika s katetama $\frac a2=3$ cm in $m=4$ cm (stranska ploskev je enakokraki trikotnik, vi\v{s}ina $m$ razpolavlja osnovnico $a$):
\[
b=\sqrt{m^2+\left(\frac a2\right)^2}=\sqrt{16+9}=5\ \text{cm}.
\]
Vsota dol\v{z}in vseh robov:
\[
o=4a+4b=4\cdot6+4\cdot5=24+20=44\ \text{cm}.
\]
\[
\boxed{o=44\ \text{cm}}
\]

\textbf{c)} Kot $\varphi$ med stransko in osnovno ploskvijo nastopa v pravokotnem trikotniku s katetama $v$ in $\frac a2=3$ ter hipotenuzo $m=4$:
\[
\cos\varphi=\frac{a/2}{m}=\frac34\ \Rightarrow\ \varphi=\arccos\frac34\approx41{,}41\degree.
\]
Ker je $0{,}41\degree\cdot60'\approx25'$, je
\[
\boxed{\varphi\approx41\degree25'}
\]
\end{resitev}

\newpage
%%%%%%%%%%%%%%%%%%%% naloga 4
\begin{resitev}[-- sto\v{z}ec]
Osni presek sto\v{z}ca je enakokrak trikotnik z osnovnico $2r$, vi\v{s}ino $v$ in kraki, ki so enaki stranici sto\v{z}ca $s$.

\textbf{a)} Plo\v{s}\v{c}ina osnega preseka:
\[
\frac{2r\cdot v}{2}=r\cdot v=60\ \Rightarrow\ r\cdot12=60\ \Rightarrow\ r=5\ \text{cm}.
\]
Stranica sto\v{z}ca (Pitagorov izrek): $s=\sqrt{r^2+v^2}=\sqrt{25+144}=13$ cm. Povr\v{s}ina pla\v{s}\v{c}a:
\[
S_{pl}=\pi rs=\pi\cdot5\cdot13=65\pi\ \text{cm}^2.
\]
\[
\boxed{S_{pl}=65\pi\ \text{cm}^2\approx204{,}20\ \text{cm}^2}
\]

\textbf{b)} Veliki krog sto\v{z}cu vri\v{s}ene krogle je hkrati vrisana kro\v{z}nica osnega preseka -- trikotnika s stranicami $13$, $13$ in $10$ cm. Njegova plo\v{s}\v{c}ina je $60\ \text{cm}^2$, polobseg pa $p=\dfrac{13+13+10}{2}=18$ cm. Iz $P=\rho\cdot p$ sledi
\[
\rho=\frac{P}{p}=\frac{60}{18}=\frac{10}{3}\ \text{cm}.
\]
\[
\boxed{\rho=\frac{10}{3}\ \text{cm}\approx3{,}33\ \text{cm}}
\]
\end{resitev}

\end{document}
